\section*{Chinesische Nudeln}
\ihead{}\chead{Personen:1}\ohead{}
\ifoot{Vorbereitungszeit:20 min}\cfoot{}\ofoot{Kochzeit:15 min}
{\Large Zutaten}
\begin{multicols}{2}
\begin{itemize}
\item \SI{125}{g} Woknudeln
\item \SI{80}{g} Hähnchenfleisch
\item etwas Öl zum Braten
\item \num{1} Ei, Größe M
\item \num{1} Handvoll Karotten, in Streifen
\item \num{1}{EL} Zwiebeln, gehackt
\item \num{1} Handvoll Sojasprossen
\item \num{1} Handvoll Lauch
\item \SI{1/4}{TL} Zucker
\item \num{1}{EL} dunkle Sojasauce
\item etwas Salz und weißer Pfeffer
\item \SI{1/2}{TL} Sesamöl
\end{itemize}
% \columnbreak
\end{multicols}
\noindent {\Large Zubereitung}\\
\\
\textit{Vorbereitung:} Die Nudeln mit warmem Wasser übergießen und ca. \SI{15}{min} ziehen lassen bzw. nach Packungsangabe kochen, bis sie gar, aber nicht zu weich sind, und anschließend abtropfen lassen.\
Das Hähnchenfleisch in kleine Stücke schneiden, in gesalzenem Wasser kochen und anschließend abtropfen lassen.\
Etwas Öl in einen heißen Wok oder eine Pfanne geben und darin das verquirlte Ei kurz braten und verrühren.\
Die gehackte Zwiebel hinzufügen und kurz weiterbraten.\
Die Karottenstreifen und den klein geschnittenen Lauch dazugeben und kurz mitbraten.\
Anschließend die Nudeln hinzufügen und mit Salz, Zucker und der dunklen Sojasauce würzen.\
Alles gut mischen und weiterbraten.\
Das Hähnchenfleisch und die Sojasprossen hinzufügen und kurz unter hoher Hitze untermischen.\
Etwas weißen Pfeffer und ca. \SI{0,5-1}{TL} Sesamöl unterrühren.\
Die gebratenen Nudeln servieren.
